\section{Exact certificates for the disjoint-support gap}
\label{family:appendix}

This appendix gives the complete rational certificate used in
\cref{family:gap}. All its entries are finite rational numbers, so
positive definiteness and the SOC comparisons reduce to integer arithmetic.

\subsection{Moments and localizing matrices}

\begin{table}[htbp]
\centering
\caption{The functional $\Lambda$ on $V$. Each listed numerator is divided
by $10000$. The exponent $(a,b,c)$ represents $x^ay^bz^c$.}
\label{family:moment-table}
\begin{tabular}{rrrr}
\toprule
Exponent & Numerator & Exponent & Numerator\\
\midrule
$(0,0,0)$ & $10000$ & $(1,0,0)$ & $4041$\\
$(0,0,1)$ & $1377$  & $(1,0,1)$ & $1025$\\
$(0,0,2)$ & $681$   & $(1,0,2)$ & $629$\\
$(0,1,0)$ & $4041$  & $(1,1,0)$ & $771$\\
$(0,1,1)$ & $1025$  & $(1,1,1)$ & $674$\\
$(0,1,2)$ & $629$   & $(1,1,2)$ & $590$\\
$(0,2,0)$ & $3392$  & $(1,2,0)$ & $770$\\
$(0,2,1)$ & $1222$  & $(1,2,1)$ & $598$\\
$(2,0,0)$ & $3392$  & $(2,1,0)$ & $770$\\
$(2,0,1)$ & $1222$  & $(2,1,1)$ & $598$\\
\bottomrule
\end{tabular}
\end{table}

For a disjoint pair $(A,B)$, let $u_{A,B}$ list $1$ followed by the
coordinates outside $A\cup B$, in increasing coordinate order. The
localizing matrix is
\[
M_{A,B}=\Lambda(w_{A,B}u_{A,B}u_{A,B}^T).
\]
The following tables give every such matrix, with each matrix multiplied
by $10000$. Their entries follow directly by expanding the factors
$1-x_i$ against \cref{family:moment-table}.
The moments are invariant under exchanging $x$ and $y$; rows marked with
multiplicity two represent the corresponding pair of matrices, with the
free coordinates exchanged. This convention preserves the displayed
matrix and the order of its basis.

\begin{proposition}
\label{family:localizing-tables}
All $27$ localizing matrices specified by \cref{family:moment-table} are
positive definite. They consist of one order-four matrix, six order-three
matrices, twelve order-two matrices, and eight scalar matrices.
\end{proposition}

\begin{proof}
The unweighted matrix, with basis $(1,x,y,z)$, is
\begin{equation}
\label{family:orderfour}
10000M_{\varnothing,\varnothing}
=\begin{pmatrix}
10000&4041&4041&1377\\
4041&3392&771&1025\\
4041&771&3392&1025\\
1377&1025&1025&681
\end{pmatrix}.
\end{equation}
Its four leading determinants, before dividing the matrix by $10000$, are
\[
10000,\quad 17590319,\quad 23512042198,\quad 45063248771.
\]

For the order-three blocks, let $u$ be either $x$ or $y$, and let $v$ be
the other coordinate. The basis for a weight involving $u$ is $(1,v,z)$;
the basis for a weight involving $z$ is $(1,x,y)$.
\begin{center}
\renewcommand{\arraystretch}{1.45}
\begin{tabular}{ccc}
\toprule
Weight & $10000M_{A,B}$ & Multiplicity\\
\midrule
$1-u$ & $\begin{pmatrix}5959&3270&352\\3270&2622&351\\352&351&52\end{pmatrix}$ & $2$\\
$u$   & $\begin{pmatrix}4041&771&1025\\771&770&674\\1025&674&629\end{pmatrix}$ & $2$\\
$1-z$ & $\begin{pmatrix}8623&3016&3016\\3016&2170&97\\3016&97&2170\end{pmatrix}$ & $1$\\
$z$   & $\begin{pmatrix}1377&1025&1025\\1025&1222&674\\1025&674&1222\end{pmatrix}$ & $1$\\
\bottomrule
\end{tabular}
\end{center}
Their leading determinants are as follows; the order-$r$ column must be
divided by $10000^r$ to obtain the determinant of the unscaled block.
\begin{center}
\begin{tabular}{rrrr}
\toprule
Weight & Order $1$ & Order $2$ & Order $3$\\
\midrule
$1-u$ & $5959$ & $4931598$ & $5442129$\\
$u$ & $4041$ & $2517129$ & $3854275$\\
$1-z$ & $8623$ & $9615654$ & $2810633517$\\
$z$ & $1377$ & $632069$ & $279229016$\\
\bottomrule
\end{tabular}
\end{center}

For the order-two blocks, the basis is $(1,u)$ in the first four rows
below, where $\{u,v\}=\{x,y\}$, and $(1,z)$ in the final three rows.
\begin{center}
\renewcommand{\arraystretch}{1.4}
\begin{tabular}{ccrr}
\toprule
Weight & $10000M_{A,B}$ & Multiplicity & Determinant numerator\\
\midrule
$(1-v)(1-z)$ & $\begin{pmatrix}5607&2919\\2919&1998\end{pmatrix}$ & $2$ & $2682225$\\
$v(1-z)$ & $\begin{pmatrix}3016&97\\97&172\end{pmatrix}$ & $2$ & $509343$\\
$(1-v)z$ & $\begin{pmatrix}352&351\\351&624\end{pmatrix}$ & $2$ & $96447$\\
$vz$ & $\begin{pmatrix}1025&674\\674&598\end{pmatrix}$ & $2$ & $158674$\\
$(1-x)(1-y)$ & $\begin{pmatrix}2689&1\\1&13\end{pmatrix}$ & $1$ & $34956$\\
$x(1-y),\ (1-x)y$ & $\begin{pmatrix}3270&351\\351&39\end{pmatrix}$ & $2$ & $4329$\\
$xy$ & $\begin{pmatrix}771&674\\674&590\end{pmatrix}$ & $1$ & $614$\\
\bottomrule
\end{tabular}
\end{center}
Each determinant numerator is divided by $10000^2$. Every first diagonal
entry is positive. Finally, the eight scalar blocks, indexed by the cube
vertex corresponding to their literal weight, are
\begin{center}
\begin{tabular}{rrrrrrr}
\toprule
Vertex & $000$ & $001$ & $010,100$ & $011,101$ & $110$ & $111$\\
\midrule
Numerator & $2688$ & $1$ & $2919$ & $351$ & $97$ & $674$\\
\bottomrule
\end{tabular}
\end{center}
Again all entries are divided by $10000$. The multiplicities give all
$27$ blocks. Positivity of every displayed leading determinant proves
the assertion by Sylvester's criterion.
\end{proof}

For reference, the minimum leading principal determinant of each order,
over every block where that order is present, is
\begin{equation}
\label{family:determinant-minima}
\frac1{10000},\qquad
\frac{307}{50000000},\qquad
\frac{154171}{40000000000},\qquad
\frac{45063248771}{10^{16}}.
\end{equation}
The tables provide a positive-definiteness certificate without solving
an SDP. The separating objective can likewise be reproduced directly:
\[
10000\Lambda(p)
=2(3392)+9(681)+6(771)-24(1025)-2(4041)+9(1377)+2500
=-250.
\]

\subsection{All permuted and complemented SOC inequalities}

For $s=(s_x,s_y,s_z)\in\{0,1\}^3$, set
$u_i=x_i$ when $s_i=0$ and $u_i=1-x_i$ when $s_i=1$.
Write $Y^s_{ij}=\Lambda(u_iu_j)$ and $t_s=\Lambda(u_xu_yu_z)$.
The entries needed for \eqref{family:ap-soc} are in the following table.
As before, each entry is divided by $10000$. Exchanging $x$ and $y$
gives the two omitted rows, $010$ and $011$.
\begin{center}
\begin{tabular}{rrrrrrrr}
\toprule
$s$ & $Y^s_{xx}$ & $Y^s_{yy}$ & $Y^s_{zz}$ & $Y^s_{xy}$ & $Y^s_{xz}$ & $Y^s_{yz}$ & $t_s$\\
\midrule
$000$ & $3392$ & $3392$ & $681$ & $771$ & $1025$ & $1025$ & $674$\\
$100$ & $5310$ & $3392$ & $681$ & $3270$ & $352$ & $1025$ & $351$\\
$001$ & $3392$ & $3392$ & $7927$ & $771$ & $3016$ & $3016$ & $97$\\
$110$ & $5310$ & $5310$ & $681$ & $2689$ & $352$ & $352$ & $1$\\
$101$ & $5310$ & $3392$ & $7927$ & $3270$ & $5607$ & $3016$ & $2919$\\
$111$ & $5310$ & $5310$ & $7927$ & $2689$ & $5607$ & $5607$ & $2688$\\
\bottomrule
\end{tabular}
\end{center}

\begin{proposition}
\label{family:ap-certificate}
The functional in \cref{family:moment-table} satisfies all permuted and
complemented inequalities in \eqref{family:ap-soc} strictly. The minimum
slacks are $2831/4000000$ for the first type and
$124813/12500000$ for the second type. All right-hand factors are strictly
positive.
\end{proposition}

\begin{proof}
For each complement pattern, form the three first-type slacks and the
six second-type slacks
\[
Y^s_{ii}Y^s_{jk}-t_s^2,\qquad
Y^s_{ii}(Y^s_{jj}+3Y^s_{jk})-(Y^s_{ij}+t_s)^2,
\]
respectively. Direct substitution from the preceding table gives the
following minima. The entries are integer numerators with denominator
$10^8$.
\begin{center}
\begin{tabular}{rrr}
\toprule
Complement pattern & Minimum first-type slack & Minimum second-type slack\\
\midrule
$000$ & $70775$ & $998504$\\
$100,010$ & $1070783$ & $3998528$\\
$001$ & $6102308$ & $35532766$\\
$110$ & $1831208$ & $8985128$\\
$101,011$ & $7494399$ & $17444574$\\
$111$ & $14090359$ & $37232454$\\
\bottomrule
\end{tabular}
\end{center}
These eight patterns cover all $24$ first-type inequalities and all $48$
second-type inequalities. Their global minima reduce to the stated
fractions. Every diagonal and cross moment in the table is strictly
positive, which supplies the nonnegative factors required for the rotated
SOC inequalities. The smallest such diagonal or cross factor is
$352/10000=22/625$.
\end{proof}

The same scalar blocks also verify the ordinary multilinear baseline.
If $\lambda_s$ denotes the scalar block at vertex $s$, then
$\lambda_s>0$ and $\sum_s\lambda_s=1$. Marginalization gives all first
and off-diagonal second moments. In particular, the off-diagonal RLT
inequalities follow by summing suitable $\lambda_s$. The two ordinary
triangle types follow from the identities
\begin{align*}
m_i+Y_{jk}-Y_{ij}-Y_{ik}
&=\Lambda\bigl(x_i(1-x_j)(1-x_k)+(1-x_i)x_jx_k\bigr),\\
1-m_x-m_y-m_z+Y_{xy}+Y_{xz}+Y_{yz}
&=\Lambda\bigl((1-x)(1-y)(1-z)+xyz\bigr).
\end{align*}
Together with the positive definite base moment matrix and the diagonal
upper-bound slacks in \cref{family:gap}, this verifies the claimed
baseline for the comparison with \citet{AnstreicherPuges2025ETRI}.

\subsection{Closedness of the PSD--nonnegative decomposition}

The order-four cone identity used in \cref{family:lmi} is an equality of
cones, rather than a closure statement. The following elementary
closedness argument also explains why dualizing $\CP_4=\DNN_4$ does
not introduce a missing limit case.

\begin{lemma}
\label{family:spn-closed}
For every $r$, the cone $\Sbb^r_++\mathcal N_r$ is closed, where
$\mathcal N_r$ is the symmetric entrywise nonnegative cone.
\end{lemma}

\begin{proof}
Suppose $P_j+N_j\to A$, with $P_j\succeq0$ and $N_j\ge0$ entrywise.
Their diagonal entries are nonnegative and bounded, since their sums
converge. The PSD inequality
$|(P_j)_{ab}|^2\le(P_j)_{aa}(P_j)_{bb}$ bounds every entry of $P_j$.
Then $N_j=(P_j+N_j)-P_j$ is bounded as well. A common convergent
subsequence yields $A=P+N$, with $P\succeq0$ and $N\ge0$ entrywise.
\end{proof}
